The course notes anchor the equal-temperament scale at $A_4=440\ \mathrm{Hz}$ and divide each octave into twelve equal semitone intervals. If $n$ is the number of semitones above or below $A_4$, the frequency is
\[
f_n=440\,2^{n/12}\ \mathrm{Hz}.
\]
The input values and the note range shown in the course-note chart are
\begin{center}\small
\begin{tabular}{@{}lrl@{}}
\toprule\noalign{}
Quantity & Value & Unit \\
\midrule\noalign{}
$A_4$ & 440 & Hz \\
$f$ & 246.94 & Hz \\
Chart range & $C_0$--$C_{10}$ & notes \\
\bottomrule
\end{tabular}
\end{center}
The chart appears on page 90 of the course notes.

For A.5.1, $C_5$ is three semitones above $A_4$. For A.5.2, the frequency-to-note relation gives the semitone offset
\[
n=12\log_2\!\left(\frac{246.94}{440}\right)=-10.0001\approx-10,
\]
so the supplied frequency corresponds to $B_3$, ten semitones below $A_4$.
\[
\begin{aligned}
f_{C_5}&=440\,2^{3/12}=523.251131\ \mathrm{Hz}\approx\boxed{523\ \mathrm{Hz}} &&\text{(A.5.1)},\\
f_{B_3}&=440\,2^{-10/12}=246.941651\ \mathrm{Hz}\approx246.94\ \mathrm{Hz},\\
246.94\ \mathrm{Hz}&\longleftrightarrow\boxed{B_3} &&\text{(A.5.2)}.
\end{aligned}
\]

The course-note chart begins at $C_0$ and ends at $C_{10}$. Applying the same equation gives its endpoints. Equal temperament itself continues by octave in both directions, so the mathematical scale has no absolute lowest or highest frequency.
\[
\begin{aligned}
f_{C_0}&=440\,2^{-57/12}=16.351598\ \mathrm{Hz}\approx\boxed{16.35\ \mathrm{Hz}\text{ at }C_0} &&\text{(A.5.3, chart minimum)},\\
f_{C_{10}}&=440\,2^{63/12}=16744.036179\ \mathrm{Hz}\approx\boxed{16.74\ \mathrm{kHz}\text{ at }C_{10}} &&\text{(A.5.5, chart maximum)}.
\end{aligned}
\]

The commonly accepted hearing range is approximately $20\ \mathrm{Hz}$ to $20\ \mathrm{kHz}$.
\[
\boxed{20\ \mathrm{Hz}\leq f\leq20\ \mathrm{kHz}}\qquad\text{(A.5.4)}
\]

Equal temperament spaces the twelve semitone steps uniformly on a logarithmic frequency scale. Each successive entry is therefore the previous frequency multiplied by the same factor, and twelve such steps double the frequency to reach the same note in the next octave.
\[
\boxed{\begin{aligned}
\text{entries per octave}&=12,\\
\frac{f_{n+1}}{f_n}&=2^{1/12}\approx1.059,\\
\left(2^{1/12}\right)^{12}&=2.
\end{aligned}}\qquad\text{(A.5.6)}
\]

\clearpage
\matlabcaption{MATLAB calculations for A.5.}
\lstinputlisting[style=hwmatlab]{matlab/A5.m}
